\chapter{Introduction}
\label{ch:intro}

\section{Purpose}

This document presents the design for the GENESYS Code Import/Export Tool in
Unified Modeling Language. It expresses the requirements agreed in the System
Requirements Specification as a structural and behavioural design, so that
implementation can proceed from a shared understanding of the system.

Each diagram is accompanied by an explanation of its elements, the design
decisions it records, and the requirements it realises. Chapter~\ref{ch:traceability}
gathers those mappings into a single matrix.

\section{System Under Design}

The tool translates bidirectionally between models held in a GENESYS repository
and source code held in an external repository. Three translations are specified:
state machine to activity diagram, activity diagram to code, and code to state
diagram. Each has its own activity diagram in Chapter~\ref{ch:activity}.

For the minimum viable product it is a terminal-based .NET application running in
parallel to GENESYS and communicating through the GENESYS API, as stated in
Section~2.3 of the requirements. Invoking it from the GENESYS interface, with the
application running as a background service driven by a GENESYS script, is a
stretch requirement rather than a committed one. The design therefore keeps the
application separable from any interface placed in front of it.

\section{Document Conventions}

Every diagram carries a figure number and a caption, is referenced from the
surrounding text, and is followed by a table explaining its elements. Where a
diagram is still provisional, the gaps are stated explicitly rather than left for
the reader to notice.

The requirements specification states its requirements as bullet lists without
identifiers. For traceability, this document assigns an identifier to each,
numbered in the order it appears:

\begin{itemize}[nosep]
	\item \req{SA-n} --- Section 3.1.1, State Diagrams to Activity Diagrams.
	\item \req{AC-n} --- Section 3.1.2, Activity Diagrams to Code.
	\item \req{CS-n} --- Section 3.1.3, Code to State Diagram.
	\item \req{FS-n} --- Section 3.1.4, Stretch Functional Requirements.
	\item \req{NFR-n} --- Section 3.2.1, Non-Functional, minimum viable product.
	\item \req{NFS-n} --- Section 3.2.2, Non-Functional, stretch.
\end{itemize}

The full list, with the text each identifier stands for, is in
Table~\ref{tab:traceability}.

\section{Diagram Index}

\begin{longtable}{L{0.20\textwidth}L{0.42\textwidth}L{0.28\textwidth}}
	\caption{Diagrams in this document}\label{tab:diagram-index}\\
	\toprule
	\textbf{Diagram} & \textbf{Models} & \textbf{Chapter} \\
	\midrule
	\endfirsthead
	\toprule
	\textbf{Diagram} & \textbf{Models} & \textbf{Chapter} \\
	\midrule
	\endhead
	Block    & System composition and boundaries         & Chapter~\ref{ch:block} \\
	Use case & Interaction between the actor and the tool & Chapter~\ref{ch:usecase} \\
	Activity & The work performed in each of the three translations, one diagram each & Chapter~\ref{ch:activity} \\
	Sequence & Interaction between participants over time & Chapter~\ref{ch:sequence} \\
	Class    & Static structure of the implementation     & Chapter~\ref{ch:class} \\
	State    & Lifecycle of a tool session                & Chapter~\ref{ch:state} \\
	\bottomrule
\end{longtable}
